\chapter{Výběr RTK GNSS/IMU systému}
\label{chap:vyber}

\section{Vymezení problému a cíl rešerše}

Cílem navrhovaného systému je rekonstruovat trojrozměrnou trajektorii vozidla
Formula Student a současně zaznamenat jeho rychlou dynamiku. Výsledkem nemá být
jen stopa vykreslená na mapě, ale měřicí záznam s dohledatelnou časovou osou,
stavem RTK řešení, odhadem nejistoty a možností opakovaného zpracování. Systém
proto musí souběžně vytvářet tři oddělené datové proudy:

\begin{enumerate}
  \item \textbf{surová data GNSS}: pseudovzdálenosti, nosné fáze, Dopplerův
        posuv, navigační zprávy, přijaté korekce, čas přijímače a diagnostika;
  \item \textbf{surová data IMU}: zrychlení, úhlové rychlosti, magnetické
        pole, teplota, stavová slova, softwarové pořadí a podklady k časování;
  \item \textbf{řešení v reálném čase}: poloha, rychlost, orientace,
        kovariance, stav RTK a diagnostické příznaky fúze.
\end{enumerate}

První dva proudy představují primární měření. Třetí proud slouží během jízdy
pro telemetrii a rychlé vyhodnocení. Po jízdě se surová data zpracují znovu,
tentokrát bez omezení na kauzální algoritmus. Lze využít PPK, dopředný a zpětný
průchod filtrem nebo vyhlazování celé trajektorie. Takové rozdělení chrání
projekt před situací, kdy chyba v prvním filtru nenávratně znehodnotí měření.

Při zpřesnění zadání jsem oddělil přesnost pohybu \textbf{vůči základnové
stanici} od přesnosti umístění celé stopy na mapě. Rozhodující je pro mě
první z nich. Trajektorii proto chci vyhodnocovat především v místním
souřadném systému ENU (východ, sever, nahoru), navázaném na nehybnou
základnu. Přibližné absolutní umístění na mapě mi pro první etapu postačuje.
Neznamená to rezignaci na přesnost měření, ale omezení nákladů na přesné
geodetické určení počátku soustavy.

Rešerše porovnává veřejně dokumentované parametry dostupných komponent. Údaje
v tabulkách pocházejí přednostně z dokumentace výrobců; stav trhu je vztažen k
období 23.~až 30.~září~2026. Volby IMU, LoRa a antén i zdůvodnění
základnové architektury jsou aktualizovány k 30.~září; data cenových
podkladů jsou rozlišena v rozpočtu. Číselné parametry jsou podmíněny konfigurací, prostředím,
firmwarem a licencí. Hodnoty z katalogu proto nejsou výsledkem kalibrace ani
garancí přesnosti celé sestavy.

\section{Proč nestačí samotné rychlé RTK GNSS}

RTK využívá nosnou fázi a korekce z referenční stanice. Při správně určených
celých nejednoznačnostech vzniká stav \emph{FIX}; při nevyřešených
nejednoznačnostech přijímač pracuje ve stavu \emph{FLOAT}. Centimetrová
přesnost uváděná výrobcem se vztahuje k vhodným podmínkám, kvalitní anténě,
krátké základně, otevřenému horizontu a právě ke stavu FIX. Neplatí automaticky
pro průjezd pod stromy, mezi budovami nebo v blízkosti uhlíkových a kovových
částí vozidla.

Vysoká výstupní frekvence rovněž neznamená, že každá poloha obsahuje stejně
silnou novou informaci. Přijímač může v různých režimech rychle publikovat
predikované řešení, zatímco nové surové observace nebo RTK aktualizace vznikají
pomaleji. U každého výrobku je proto nutné rozlišovat:

\begin{itemize}
  \item frekvenci sledování a surových observací;
  \item frekvenci výpočtu absolutní polohy;
  \item frekvenci skutečného RTK řešení;
  \item latenci a čas, ke kterému se publikovaná poloha vztahuje.
\end{itemize}

IMU naopak změří dynamiku s malou latencí a vysokou frekvencí, ale integrací
jejího biasu vzniká rostoucí chyba rychlosti, polohy a orientace. GNSS a IMU se
proto doplňují: GNSS dlouhodobě ukotvuje stav v souřadném systému základny
s možností převodu do globálních souřadnic,
zatímco IMU interpoluje pohyb mezi GNSS epochami a pomáhá překlenout krátké
výpadky. Přesnost kombinace určuje nejen kvalita senzorů, ale především společný
čas, kalibrace orientace a přesně změřená páka mezi fázovým středem antény a
měřicím bodem IMU.

\section{Návrhové požadavky}

Požadavky v tabulce~\ref{tab:pozadavky} jsou navrženy pro výzkumný logger.
Systém není v první etapě součástí bezpečnostní funkce vozidla a jeho chyba
nesmí ovlivnit řízení, odpojovací smyčku ani jiné povinné systémy. Instalace
však musí splnit aktuální pravidla konkrétního závodu. Pravidla Formula Student
Germany 2026 požadují dobrou inženýrskou praxi a u autonomního vozidla také
dokumentaci použitých percepčních senzorů a jejich souladu s místní
legislativou.
\cite{fsg-rules-2026}

\begin{table}[htbp]
  \centering
  \small
  \caption{Návrhové požadavky na roverový měřicí systém}
  \label{tab:pozadavky}
  \begin{tabularx}{\textwidth}{@{}L{0.29\textwidth}L{0.30\textwidth}Y@{}}
    \toprule
    Vlastnost & Minimální akceptovatelná hodnota & Cílová hodnota \\
    \midrule
    Relativní RTK poloha vůči základně, FIX
      & vodorovně do \SI{2}{cm}, svisle do \SI{3}{cm}
      & ověřená trajektorie s kovariancí a příznakem kvality \\
    GNSS observace / RTK řešení
      & \SI{20}{Hz} / \SI{20}{Hz}
      & \SI{100}{Hz} / alespoň \SI{50}{Hz} \\
    Navigační IMU
      & nejméně \SI{400}{Hz}, \SI{\pm8}{g}, \SI{\pm1000}{\degree\per\second}
      & nominálně \SI{562.5}{Hz}, \SI{\pm16}{g}, \SI{\pm2000}{\degree\per\second} \\
    Časová vazba
      & GNSS PPS, události GPIO a odhad nejistoty času
      & experimentálně ověřená chyba epoch; pracovní cíl pod \SI{100}{\micro\second} \\
    Fúzovaný výstup
      & \SI{100}{Hz}
      & \SIrange{200}{500}{Hz} včetně orientace a kovariance \\
    Záznam
      & tři oddělené binární proudy
      & kontrolní součty, metadata, konfigurace a události \\
    Odolnost
      & bez ztráty dat při výpadku korekcí
      & bezpečné ukončení při poklesu napájení a automatické zotavení \\
    Provozní teplota sestavy
      & \SIrange{0}{60}{\degreeCelsius}
      & návrh a ověření nejméně pro \SIrange{-10}{70}{\degreeCelsius} \\
    \bottomrule
  \end{tabularx}
\end{table}

Hodnoty přesnosti v tabulce jsou pracovní cíle pro místní trajektorii,
nikoli garance absolutní polohy na mapě. Splnění posoudím porovnáním
s nezávislou referencí a uvedu podmínky i statistiku chyby; samotný příznak
FIX ani kovariance přijímače takový důkaz nepředstavují.

Podle upřesněného zadání aktivuje ochrana nárazu vozidla vypnutí přibližně při
\SI{3}{g}. Rozsah \SI{\pm16}{g} proto považuji za přiměřený první návrh.
Práh ochrany však není horní mez okamžitého zrychlení v místě IMU. Lokální vibrace a mechanický ráz v místě snímače mohou
být krátkodobě vyšší než zrychlení těžiště, proto musí logger stále ukládat
maximum každé osy a příznak saturace. Samostatný \SI{\pm200}{g} kanál však pro
první prototyp není nutný. Plastový kryt se nesmí považovat za kalibrovaný
mechanický filtr: jeho tuhost, vůle a vlastní rezonance je třeba ověřit z
naměřeného spektra.

Při \SI{100}{km\per\hour}, tedy přibližně \SI{27.78}{m\per\second}, ujede
vozidlo mezi vzorky \SI{50}{Hz} asi \SI{0.556}{m}; při \SI{500}{Hz} asi
\SI{5.56}{cm}. Zvolené nominální ODR \SI{562.5}{Hz} nové IMU odpovídá
přibližně \SI{4.94}{cm}. Obecně platí
\begin{equation}
  \Delta s = \frac{v}{f_\mathrm{s}},
\end{equation}
kde $v$ je rychlost vozidla a $f_\mathrm{s}$ vzorkovací frekvence. Tento
výpočet popisuje prostorový rozestup inerciálních vzorků,
nikoli přesnost polohy: výslednou přesnost určuje především GNSS, časová
synchronizace, bias IMU a kvalita fúze. Pro sledování trajektorie a pohybu karoserie předpokládám, že přibližně
\SI{500}{Hz} postačí. Tento předpoklad ověřím proti naměřené šířce pásma;
pro analýzu vysokofrekvenčních vibrací může být nutné jiné nastavení.

\section{Průzkum trhu RTK GNSS přijímačů}

\subsection{Hodnoticí kritéria}

Pro rover jsou rozhodující vícefrekvenční příjem, přístup k surovým
observacím, nejméně jeden konfigurovatelný PPS výstup, nízká latence, příjem
RTCM~3, robustnost proti multipath a rušení, otevřená dokumentace a rozhraní
vhodná pro záznam. Samotný počet kanálů není vhodným kritériem, pokud výrobce
nedokládá chování přijímače v dynamice.

\subsection{Porovnání kandidátů}

\begin{table}[htbp]
  \centering
  \scriptsize
  \caption{Srovnání kandidátů na roverový GNSS přijímač}
  \label{tab:gnss}
  \begin{tabularx}{\textwidth}{@{}L{0.18\textwidth}L{0.15\textwidth}L{0.18\textwidth}Y@{}}
    \toprule
    Přijímač & Frekvence & Výrobní údaj RTK & Posouzení pro projekt \\
    \midrule
    Septentrio mosaic-X5 / mosaic-go X5
      & poloha až \SI{100}{Hz}; sada uvádí měření \SI{100}{Hz} a RTK \SI{50}{Hz}
      & cm třída, rychlost \SI{3}{cm\per\second}; PPS a dva event vstupy
      & Preferovaná varianta. SBF raw, USB, sériová rozhraní, SD a kvalitní
        diagnostika rušení. Před nákupem potvrdit konkrétní RTK režim.
        \cite{septentrio-mosaic,septentrio-mosaic-go,septentrio-hw-manual} \\
    Septentrio mosaic-G5 P6
      & poloha i raw až \SI{100}{Hz}
      & vodorovně \SI{0.6}{cm}+\SI{0.5}{ppm}, svisle \SI{1}{cm}+\SI{1}{ppm}
      & Novější, perspektivní modul. Vhodný jako náhrada X5, pokud dodavatel
        nabídne vývojovou desku, termín a cenu odpovídající projektu.
        \cite{septentrio-g5-p6} \\
    Unicore UM980
      & poloha a vybraná raw data až \SI{50}{Hz}
      & vodorovně \SI{0.8}{cm}+\SI{1}{ppm}, svisle \SI{1.5}{cm}+\SI{1}{ppm}
      & Velmi dobrý poměr výkon/cena; PPS \SI{20}{ns}. Vyžaduje pečlivý výběr
        nosné desky a ověření tepelného i napájecího návrhu.
        \cite{unicore-um980,unicore-protocol} \\
    u-blox ZED-F9P
      & RTK až \SI{20}{Hz} podle varianty a konfigurace
      & přibližně \SI{1}{cm}+\SI{1}{ppm}
      & Výborný pro nízkonákladovou základnu a první laboratorní prototyp;
        nesplní cílových \SI{50}{Hz} roverového RTK.
        \cite{ublox-f9p} \\
    NovAtel OEM7720
      & poloha až \SI{100}{Hz}
      & profesionální RTK/PPP ekosystém
      & Technicky silná varianta se dvěma anténními vstupy a přesným časem,
        avšak s vyšší pořizovací cenou a integračními náklady.
        \cite{novatel-oem7720} \\
    Septentrio mosaic-H
      & poloha až \SI{50}{Hz}
      & heading \SI{0.15}{\degree} při základně \SI{1}{m}
      & Dvě antény určují azimut i v klidu. Přínosné pro autonomní režim,
        ale na malé formuli přidává druhou anténu, kabeláž a přesnou základnu.
        \cite{septentrio-mosaic-h} \\
    \bottomrule
  \end{tabularx}
\end{table}

U mosaic-X5 existuje zdánlivý rozpor mezi produktovou stránkou modulu, která
uvádí maximální polohový výstup \SI{100}{Hz}, a stránkou vývojové sady
mosaic-go X5, která uvádí \SI{100}{Hz} pro měření a \SI{50}{Hz} pro RTK.
Hardware manuál navíc obsahuje příkon pro konfiguraci RTK \SI{100}{Hz}.
\cite{septentrio-mosaic,septentrio-mosaic-go,septentrio-hw-manual}
Nejde bezpečně předpokládat, že libovolná zakoupená deska poskytne
\SI{100}{Hz} centimetrové řešení se všemi konstelacemi a zprávami. Do poptávky
proto musí být uveden požadovaný firmware, výstupní zprávy, počet konstelací,
režim korekcí a požadavek na písemné potvrzení frekvence.

\subsection{Vícekriteriální posouzení GNSS}

Tabulka~\ref{tab:gnss-score} je inženýrské posouzení pro tuto konkrétní práci,
nikoli obecný žebříček výrobků. Hodnota 5 znamená nejlepší shodu s požadavkem.
Cena je hodnocena podle očekávané třídy kompletní vývojové sestavy, protože
samotný LGA modul bez nosné desky není srovnatelný s hotovým přijímačem.

\begin{table}[htbp]
  \centering
  \small
  \caption{Bodové hodnocení GNSS kandidátů pro rover (1 až 5)}
  \label{tab:gnss-score}
  \begin{tabularx}{\textwidth}{@{}Y*{5}{>{\centering\arraybackslash}p{0.10\textwidth}}@{}}
    \toprule
    Kritérium / váha & X5 & G5 P6 & UM980 & F9P & \shortstack{OEM\\7720} \\
    \midrule
    RTK a raw frekvence / 25\,\% & 4 & 5 & 4 & 2 & 5 \\
    Časování a raw data / 20\,\% & 5 & 5 & 4 & 4 & 5 \\
    Robustnost RF / 20\,\% & 5 & 5 & 3 & 3 & 5 \\
    Integrace prototypu / 15\,\% & 5 & 3 & 3 & 5 & 3 \\
    Nákladovost / 10\,\% & 3 & 2 & 5 & 5 & 1 \\
    Podpora a životnost / 10\,\% & 5 & 4 & 3 & 5 & 5 \\
    \midrule
    Vážené skóre / 5 & 4,55 & 4,30 & 3,65 & 3,65 & 4,30 \\
    \bottomrule
  \end{tabularx}
\end{table}

Frekvenční kritérium zahrnuje požadavek nejméně \SI{50}{Hz} RTK. Časování
zohledňuje PPS, event vstupy, dostupnost raw formátu a dokumentaci; robustnost
RF potom zejména multipath, odolnost proti rušení a diagnostiku. Nákladovost je
relativní odhad bez závazné nabídky a výsledky jsou zaokrouhleny na dvě místa.

\subsection{Volba přijímače a nosné desky}

Pro rover volím \textbf{Septentrio mosaic-X5 na desce ArduSimple
simpleRTK3B Pro}. Původně jsem kvůli ceně uvažoval i o variantě InCase PIN
od GNSS.store za 499,99~EUR bez DPH. Aktuální nákupní varianta simpleRTK3B Pro
za 569~EUR bez DPH umožní spojit rover, základnu i obě antény do jedné
objednávky u evropského dodavatele. Rozdíl neposuzuji pouze podle ceny desky,
ale i podle dostupnosti, podpory, dopravy a vyvedených rozhraní.
\cite{gnssstore-x5-incase,ardusimple-x5}

V rozpočtu zůstává pouze jeden X5 a jedna levnější přenosná základna ZED-F9P.
Základna nepotřebuje rychlý polohový výstup, ale kompatibilní RTCM a současný
záznam observací. Druhou stacionární stanici na FM TUL neobjednávám. G5 P6
ponechávám jako novější alternativu pro budoucí návrh; UM980 jako možnost
při nutnosti významně snížit cenu. Nejde o současně nakupované přijímače.

Jednoanténní X5 neurčuje přímo natočení karoserie v klidu. Za jízdy lze
využít směr vektoru rychlosti, ten však při smyku není totožný s podélnou osou
vozu. Tuto odlišnost musí respektovat filtr. Magnetometr nové IMU může pomoci
s orientací, ale nenahrazuje dvouanténní GNSS heading ani kalibraci.


\subsection{Co má být vlastním přínosem práce}

Při porovnávání hotových přijímačů a vývojových desek jsem si položil
otázku, zda musím pro smysluplnou práci navrhovat také vlastní GNSS
hardware. Za hlavní přínos však považuji časově provázaný sběr, vlastní
rádiový firmware, správu základny, fúzi a experimentální ověření.
Návrh RF přijímače nebo jeho nosné desky není podmínkou tohoto cíle.
Hotový přijímač může být stejně legitimním základem práce.

Vývojovou desku volím kvůli přístupu k rozhraním, měřicím bodům a vlastnímu
mechanickému uspořádání, nikoli proto, že by samotné otevřené provedení
automaticky měřilo přesněji. U základny srovnávám zejména dvě různá
provedení stejného přijímače ZED-F9P; tuto úvahu rozvádím v
oddílu~\ref{sec:base-board}.

\section{Průzkum trhu IMU}

\subsection{Kritéria výběru}

Počet os ani nejvyšší vzorkovací frekvence samy o sobě nevystihují vhodnost
IMU. Při výběru posuzuji měřicí rozsah, šířku pásma a latenci filtrů, šum,
teplotní závislost a stabilitu biasu. Pro záznam je důležité také SPI,
přerušovací výstup, FIFO a možnost odhadnout čas skutečného měření. Přebuzení
senzoru vede ke ztrátě amplitudy, kterou pozdější zpracování neobnoví.

Devět os znamená tři osy akcelerometru, tři gyroskopu a tři magnetometru.
Nejde o devět nezávislých měření polohy ani o záruku kvalitnější inerciální
navigace. Přidaný magnetometr dává jiný druh informace a současně přináší
nové požadavky na umístění, kalibraci a detekci rušení.

\subsection{Porovnání kandidátů}

\begin{table}[htbp]
  \centering
  \scriptsize
  \caption{IMU posouzené pro první prototyp; rozsahy jsou katalogová maxima}
  \label{tab:imu}
  \begin{tabularx}{\textwidth}{@{}L{0.22\textwidth}L{0.24\textwidth}Y@{}}
    \toprule
    Senzor & Rozsah / rozhraní & Posouzení pro tuto práci \\
    \midrule
    ST ISM330DHCX
      & \SI{\pm16}{g}, \SI{\pm4000}{\degree\per\second}; SPI, FIFO, INT
      & Původní šestiosá volba. Nízký šum a vhodná časová rozhraní, bez kompasu.
        \cite{st-ism330dhcx,st-ism330dhcx-ds} \\
    TDK ICM-20948 / Adafruit
      & \SI{\pm16}{g}, \SI{\pm2000}{\degree\per\second}; SPI, FIFO, INT
      & Současná volba. Magnetometr navíc, dostupná deska u RPiShopu;
        vyšší šum a čip ve stavu EOL.
        \cite{tdk-icm20948,adafruit-icm20948,tdk-icm20948-status} \\
    ST LSM6DSO32X
      & \SI{\pm32}{g}, \SI{\pm2000}{\degree\per\second}; SPI
      & Varianta při prokázané nedostatečné rezervě akcelerometru.
        \cite{st-lsm6dso32x} \\
    Bosch BMI088
      & \SI{\pm24}{g}, \SI{\pm2000}{\degree\per\second}; SPI
      & Šestiosá alternativa orientovaná na dynamiku a vibrace.
        \cite{bosch-bmi088} \\
    TDK ICM-42688-P
      & \SI{\pm16}{g}, \SI{\pm2000}{\degree\per\second}; SPI, externí hodiny
      & Alternativa pro pozdější vlastní desku; nemá magnetometr.
        \cite{tdk-icm42688p} \\
    ADIS16507-3 / SBG Pulse-40
      & až \SI{\pm40}{g}; výstup řádově \SI{2}{kHz}
      & Vyšší kalibrovaná třída. Do prvního rozpočtu ji nezařazuji.
        \cite{adis16507,sbg-pulse40} \\
    \bottomrule
  \end{tabularx}
\end{table}

\subsection{Vývoj volby od ISM330DHCX k ICM-20948}

V první variantě jsem počítal s \textbf{ISM330DHCX na desce SparkFun
SEN-19764}. Pro plánovaný sběr kolem \SI{500}{Hz} nabízela dostatečnou
frekvenční rezervu, rozsah a dokumentaci. Při snaze nakupovat od běžných
evropských distributorů jsem následně uvažoval o vývojové desce ST
STEVAL-MKI210V2K se stejným senzorem od Farnellu. Změna nosné desky tehdy
neměnila samotnou šestiosou koncepci.\cite{sparkfun-ism330dhcx,farnell-ism330dhcx}

Při dalším výběru jsem zohlednil možnost využít také kompas. Magnetometr může
dodat doplňkovou informaci o orientaci například při inicializaci nebo při
malé rychlosti, kdy směr z GNSS rychlosti není dostatečně stabilní. Proto
jsem pro první prototyp zvolil \textbf{Adafruit ICM-20948 se STEMMA QT/Qwiic},
kód RPiShopu ADA076. Při kontrole 30.~9.~2026 byla deska nabízena za
449~Kč s DPH a uveden byl jeden skladový kus. Jde o snímek dostupnosti,
nikoli o záruku pozdějšího dodání.\cite{rpishop-icm20948}

Rozhodující je pro mě souběh přidaného magnetometru, dostupnosti hotového
modulu a možnosti koupit jej společně s Raspberry Pi a LoRa sadami. Nákup
tak zůstává u dvou dodavatelů. Netvrdím, že ISM330DHCX byla technicky
nevhodná nebo obecně nedostupná; zůstává smysluplnou alternativou, pokud
nová IMU nesplní požadavky při měření.

Změna má i nevýhody. Typická katalogová hustota šumu gyroskopu je u
ICM-20948 přibližně 0,015, u ISM330DHCX 0,005
$(^{\circ}/\mathrm{s})/\sqrt{\mathrm{Hz}}$ za podmínek příslušných
datasheetů.\cite{tdk-icm20948,st-ism330dhcx-ds}
Nejde o přímý trojnásobek výsledné chyby trajektorie, ale o důvod zahrnout
šum a drift do validace. Navíc TDK uvádí ICM-20948 ve stavu
\textbf{EOL (End of Life)}.\cite{tdk-icm20948-status}
Skladový modul je použitelný pro tento prototyp; pro opakovanou výrobu bych
musel řešit náhradu. Současný výběr tedy není dlouhodobou platformovou
garancí.

\subsection{Vzorkování a využití kompasu}

Pro první měření volím pevné rozsahy \SI{\pm16}{g} a
\SI{\pm2000}{\degree\per\second}. Akcelerometr a gyroskop chci provozovat
v režimu se zapnutým digitálním filtrem a děličem nastaveným na 1.
Tabulky konfigurace datasheetu pro tento režim uvádějí vztah
\begin{equation}
  f_\mathrm{ODR}=\frac{1125}{1+d}\ \mathrm{Hz},
  \qquad d=1,\quad f_\mathrm{ODR}=\SI{562.5}{Hz}.
\end{equation}
Magnetometr AK09916 poskytuje nejvýše \SI{100}{Hz}.
\cite{tdk-icm20948,adafruit-icm20948}
V textu proto rozlišuji nominální frekvenci a skutečně změřenou periodu.
Zaokrouhlené označení \enquote{IMU kolem 500~Hz} neznamená, že všech devět
os dodává nové měření přesně pětsetkrát za sekundu.

Šířku pásma digitálních filtrů zvolím podle naměřených vibrací a požadované
dynamiky. Jejich zpoždění musím zahrnout do časového modelu. Vyšší počet
výstupních vzorků nemůže nahradit chybějící fyzikální šířku pásma.
Původní nastavení \SI{833}{Hz} patřilo k ISM330DHCX; na nový senzor jej
nepřenáším.

Kompas budu používat podmíněně. Kovové části způsobují stálé i směrově závislé
zkreslení pole (hard-iron a soft-iron), zatímco proudy ve vozidle mohou rušení
měnit za jízdy. Naplánuji proto kalibraci namontovaného modulu a porovnání při
vypnuté a zapnuté výkonové elektronice. Měření s neobvyklou velikostí pole
nebo velkou inovací filtru označím a do korekce orientace nezařadím.
Magnetický sever je navíc nutné odlišit od zeměpisného severu.

V hlavním záznamu uchovám i odmítnutá magnetická data a důvod odmítnutí.
Filtr musí fungovat také bez magnetometru, aby jeho rušení nepoškodilo
GNSS/IMU řešení. Přínos kompasu vyhodnotím porovnáním téže jízdy se zapnutými
a vypnutými magnetickými aktualizacemi; předem jej nepovažuji za prokázaný.

\subsection{Rezerva rozsahu a kontrola saturace}

Logger bude po blocích vyhodnocovat maxima jednotlivých os a počet vzorků
nad 80\,\% plného rozsahu. Pro akcelerometr jde při zvolené konfiguraci o
\SI{12.8}{g}. Hodnoty blízké mezním digitálním kódům budou označeny jako
podezření na saturaci. Nárazová ochrana kolem \SI{3}{g} ani 3D tištěný kryt
nevylučují lokální špičku, gravitační složku či rezonanci držáku.

Pokud měření ukáže nedostatečnou rezervu, zvážím LSM6DSO32X s rozsahem
\SI{\pm32}{g}. Zachoval bych koncepci přímého sběru, ale musel bych změnit
ovladač, konfiguraci, časový model a případně zapojení; nejde o prostou
výměnu bez dalších úprav. Vyšší průmyslovou IMU bych pořizoval až při
doloženém problému se stabilitou biasu nebo driftem.

\section{Výpočetní a komunikační architektura roveru}

\subsection{Role Raspberry Pi 5}

Výpočetním uzlem bude \textbf{Raspberry Pi 5 se 2~GB RAM a 64bitovým
Raspberry Pi OS}. Čtyřjádrový procesor, USB a PCIe poskytují vhodný základ
pro logger a online filtr; skutečnou výpočetní rezervu ověřím při současném
záznamu, fúzi a rádiovém přenosu.\cite{raspberry-pi5}
Počítám s Active Coolerem a vlastním NVMe SSD připojeným přes oficiální
\textbf{M.2 HAT+}. Aktuální nákupní cena HAT+ je 349~Kč s DPH, nikoli původní
orientační částka 400~Kč. HAT+ podporuje formáty SSD 2230 a 2242; vlastní
disk formátu 2280 by vyžadoval jiný adaptér.\cite{rpishop-m2hat}

\subsection{Přímý sběr bez pomocného časovacího mikrokontroléru}

Desku Adafruit připojím přes pájecí vývody SPI k Raspberry Pi a její INT
na samostatný GPIO. STEMMA QT/Qwiic konektor vyvádí I2C, nikoli SPI.
Desku budu napájet z \SI{3.3}{V}; její převodníky úrovní umožňují připojení
SPI a INT k odpovídající logice Raspberry Pi. Pin FS/FSYNC je výjimka:
pracuje na \SI{1.8}{V} a nesmí na něj být přímo přiveden signál
\SI{3.3}{V}.\cite{adafruit-icm20948,adafruit-icm20948-pcb}

Pro sběr chci využít FIFO a přerušení data-ready. Nejprve ověřím režim,
ve kterém jednotlivé hrany odpovídají požadovaným novým inerciálním vzorkům.
Obsluha přečte stavové registry, počet bajtů a dostupná data v bloku.
Magnetometr při SPI komunikaci zpřístupní vnitřní pomocné I2C rozhraní
senzoru; jeho nové hodnoty musím odlišit od opakovaně čteného posledního
měření.\cite{tdk-icm20948}

Na ICM-20948 nelze bez úpravy přenést FIFO tagy, konfiguraci watermarku ani
interní časové značky z původního návrhu s ISM330DHCX. V první implementaci
nepředpokládám samostatný absolutní hardwarový timestamp u každého raw
vzorku. Časovou osu sestavím z pořadí vzorků, skutečně změřené periody a
zachycených událostí. Přetečení, reset nebo nejednoznačnost přiřazení
označím jako přerušení časové kontinuity.

Hrany INT i GNSS PPS lze zaznamenávat pomocí GPIO character-device API v2.
Kernel k události připojuje časovou značku a pořadové číslo.
\cite{linux-gpio-chardev}
To omezuje vliv opožděného probuzení uživatelského procesu, ale standardní
Linux ani běžný GPIO vstup negarantují bezchybný hardwarový záchyt okamžiku
hrany. FIFO chrání před krátkým zpožděním čtení; samo neřeší nejistotu epochy.
Přímé připojení je proto návrhová volba, kterou ověřím zátěžovým testem,
nikoli již prokázané deterministické řešení.

\subsection{Navržené propojení}

mosaic-X5 bude předávat data po USB, zatímco PPS povede samostatně na GPIO.
Před zapojením zkontroluji vyvedení, polaritu a napěťovou úroveň konkrétní
desky simpleRTK3B Pro. Raspberry Pi nepřipouští libovolné napěťové úrovně na
GPIO. LoRa jednotku připojím datovým USB-C kabelem; ESP32 uvnitř sady obstará
rádiové rámcování. Nepřebírá úlohu časovacího mikrokontroléru pro IMU.

\begin{figure}[htbp]
  \centering
  \begin{tikzpicture}[
      block/.style={draw,rounded corners,align=center,text width=36mm,
        minimum height=13mm,inner sep=3pt,font=\small,fill=tul!8},
      data/.style={-{Latex[length=2mm]},thick},
      timing/.style={-{Latex[length=2mm]},thick,dashed},
      lab/.style={font=\scriptsize,fill=white,inner sep=2pt}
    ]
    \node[block] (base) at (0,0) {ZED-F9P\\ANN-MB2-00};
    \node[block] (tx) at (5,0) {XIAO \(+\) Wio\\správa \(+\) raw log*};
    \node[block] (rx) at (10,0) {XIAO \(+\) Wio\\rover};
    \node[block] (gnss) at (0,-3) {mosaic-X5\\ANN-MB2-00};
    \node[block] (pi) at (5,-3) {Raspberry Pi 5\\Raspberry Pi OS};
    \node[block,text width=42mm] (ssd) at (10,-3) {\mbox{SSD \(+\) M.2 HAT+}\\\mbox{\footnotesize 3 proudy \(+\) metadata}};
    \node[block] (imu) at (5,-5.8) {Adafruit ICM-20948\\akcelerometr, gyro, kompas};
    \draw[data,<->] (base) -- node[lab,above=8mm]{UART / RTCM \(+\) UBX} (tx);
    \draw[data] (tx) -- node[lab,above=8mm]{LoRa 868 MHz} (rx);
    \draw[data] (rx.south) -- ++(0,-.7) -| node[lab,pos=.7,right]{USB / RTCM} (pi.north);
    \draw[data,<->] (gnss.east) -- node[lab,above=8mm]{USB / SBF \(+\) RTCM} (pi.west);
    \draw[timing] (gnss.south) -- ++(0,-.6) -| node[lab,pos=.3,below]{PPS / GPIO} ([xshift=-10mm]pi.south);
    \draw[data] (imu.north) -- node[lab,right]{SPI \(+\) INT} (pi.south);
    \draw[data] (pi) -- node[lab,above]{PCIe} (ssd);
  \end{tikzpicture}
  \caption{Cílová architektura se samostatnou základnou. Hvězdička označuje
  raw log vyžadující doplnění lokálního úložiště; při vývoji jej zajistí
  notebook přes USB. Ze čtyř LoRa sad jsou ve spoji dvě. Čárkovaně je PPS.}
  \label{fig:architektura}
\end{figure}

\section{Časová synchronizace a datové formáty}

\subsection{Jednotná časová osa}

PPS označuje hranici sekundy a zpráva GNSS jí přiřadí číslo týdne a sekundy.
Z opakovaných PPS událostí odhadnu převod mezi monotónními hodinami Linuxu
a GNSS časem. V této základně zaznamenám také události IMU. Z pořadí měření
a naměřené periody vytvořím odhad epoch; do modelu zahrnu zpoždění digitálních
filtrů i jejich případnou odlišnost pro akcelerometr a gyroskop.

Každý inerciální záznam bude obsahovat surová slova, softwarové pořadové
číslo, odhad času a příznaky kvality. Zachovám i původní timestampy událostí,
stavy FIFO a údaje nutné k pozdější rekonstrukci. Softwarový čítač sám
neprokazuje, že senzor žádný vzorek neztratil; musí se porovnat s periodou,
počtem dat a indikací přetečení. Po ztrátě vazby nelze chybějící časovou
informaci potichu doplnit jako jisté měření.

Čas příchodu USB paketu ani okamžik ukončení SPI čtení nejsou automaticky
časem fyzikálního měření. Přesnost zvolené metody ověřím osciloskopem nebo
logickým analyzátorem. Jako pracovní cíl používám chybu časového přiřazení
pod \SI{100}{\micro\second}; nejde o specifikaci Raspberry Pi ani o výsledek
již provedeného testu. Při rychlosti \SI{100}{km\per\hour} odpovídá časová
chyba \SI{1}{ms} posunu přibližně \SI{2.78}{cm}.

Pokud základní metoda nevyhoví, prověřím pomocnou vazbu přes FSYNC.
Nejdříve je nutné vyřešit převod na \SI{1.8}{V}, způsob zachycení události
a význam příslušných registrů. Pouhé přivedení PPS na tento pin nepovažuji
za hotovou synchronizaci. Převodník ani tato rozšiřující varianta zatím
nejsou součástí oceněné základní elektroniky.

\subsection{Tři proudy a metadata}

Každá jízda dostane identifikátor a uloží:
\begin{description}
  \item[\texttt{gnss.sbf}] nativní observace a navigační i diagnostické bloky X5;
  \item[\texttt{imu.bin}] typované raw záznamy akcelerometru a gyroskopu,
    nové magnetické vzorky, teplotu a podklady k časování;
  \item[\texttt{nav.bin}] online polohu, rychlost, orientaci, kovarianci
    a diagnostiku filtru;
  \item[\texttt{events.jsonl}] PPS, výpadky, přetečení, saturaci a stav napájení;
  \item[\texttt{manifest.yaml}] verze firmwaru a dekodéru, nastavení filtrů,
    rozsahy, skutečné periody, kalibraci a geometrii montáže;
    dále identifikátor základny, její použité souřadnice, způsob jejich
    určení, výšku antény a definici místního souřadného systému.
\end{description}
Události a manifest jsou doprovodná metadata, nikoli další hlavní měřicí
proud. Magnetometr zůstává součástí \texttt{imu.bin}; jeho rozdílný typ a
příznak nového vzorku zabrání vydávání opakovaných hodnot za měření při
\SI{562.5}{Hz}. Přijaté RTCM rámce zachovám v diagnostickém záznamu korekcí.
Současně se musí ukládat raw data základny potřebná pro PPK.

Primární soubory budou binární, verzované a blokově chráněné kontrolními
součty. CSV použiji pro export. Konfigurace a kalibrace se uloží pro každou
relaci tak, aby pozdější změna ovladače nezměnila význam již pořízených dat.

\subsection{Orientační datový objem}

Při navrženém inerciálním záznamu 64~B a nominálních \SI{562.5}{Hz} vzniká
36~kB/s, tedy 129,6~MB za hodinu. Samostatný typ magnetického záznamu o 32~B
při \SI{100}{Hz} přidá nejvýše 3,2~kB/s neboli 11,52~MB za hodinu.
Celkem jde přibližně o 0,141~GB/h v desetinných jednotkách, bez GNSS,
navigačního výstupu a událostí. Velikosti rámců jsou návrhové předpoklady,
nikoli velikosti interních FIFO paketů senzoru.

Celkový objem stanovím podle skutečně aktivovaných SBF bloků. SSD má
poskytnout dostatečnou kapacitu i zapisovací rezervu. Logger bude uzavírat
a synchronizovat bloky průběžně; celou jízdu nesmí držet pouze v RAM.

\section{Fúze v reálném čase a postprocessing}

V první etapě navrhuji volně vázaný error-state Kalmanův filtr (ESKF).
Stav zahrne polohu, rychlost, orientaci a biasy akcelerometru a gyroskopu.
Každé nové inerciální měření použiji k propagaci. GNSS poloha a rychlost
provedou korekci v odpovídající epoše, nikoli až v době přijetí zprávy.
Součástí modelu bude páka mezi IMU a anténou (lever arm). Magnetometr přidám
jako samostatnou pomalejší a podmíněnou aktualizaci.

Navržená publikační frekvence kombinovaného řešení je \SIrange{200}{500}{Hz}.
Je odlišná od frekvence raw IMU i GNSS. Po jízdě mohu nativní inerciální
záznam převzorkovat na přesných \SI{500}{Hz}, ale až po vhodné filtraci a
s použitím rekonstruované časové osy. Převzorkování nevytváří další nezávislá
absolutní měření polohy.

Pro postprocessing počítám s PPK z observací základny a roveru a následným
dopředným a zpětným vyhlazením trajektorie. ESKF lze doplnit
Rauchovým--Tungovým--Striebelovým vyhlazovačem; přínos této kombinace pro
GNSS/IMU lokalizaci popisuje i publikovaná studie.\cite{wang-esfk-rts}
RTKLIB je použitelný jako otevřená referenční cesta po převodu do RINEX,
ale je nutné ověřit podporu konkrétních signálů a porovnat výstup s řešením
přijímače.\cite{takasu-rtklib,rtklib-site}
Těsně vázaná integrace jednotlivých GNSS observací je případným dalším
rozšířením, nikoli nutnou podmínkou prvního prototypu.

Výstup zachová stav FIX/FLOAT/SINGLE, stáří korekcí, kovarianci, inovace a
příznaky odmítnutých měření. Vyhlazená trajektorie není automaticky
referenční pravdou: její přesnost musí být ověřena nezávislým měřením
a nesmí být zaměňována s hustotou výstupních bodů.

\section{Základnová stanice a přenos korekcí}

\subsection{Jediná přenosná základna}

Pro RTK potřebuje rover včasná referenční data ze základny nebo ze síťové
služby. V tomto návrhu je bude poskytovat \textbf{simpleRTK2B Budget s
u-blox ZED-F9P} a anténou \textbf{u-blox ANN-MB2-00}. Základna
zůstane během jízdy nehybná. Stejnou sadu lze testovat na univerzitě
a odvézt na závody. Trvalá stanice na FM TUL ani integrace do LiANE
nejsou podmínkou návrhu.

\subsection{Relativní přesnost a umístění na mapě}

ZED-F9P umožňuje určit souřadnice základny režimem survey-in nebo zadat
pevnou polohu. Chyba zadané polohy se přenáší do absolutních souřadnic
roveru; survey-in sám nezaručuje centimetrové geodetické zaměření.
\cite{ublox-f9p-integration}
Pro můj cíl však není nutné v první etapě pořizovat takové zaměření jen
kvůli mapovému podkladu. Při krátké základně se chyba jejího absolutního
umístění projeví převážně společným posunem stopy. To neznamená, že každá
relativní vzdálenost bude automaticky centimetrově přesná: nadále záleží
na observacích, vícecestném šíření, geometrii i správném FIX.

Pro měřicí relaci zafixuji použité souřadnice a mechanickou polohu antény.
Uložím definici počátku ENU, orientaci os a referenční výšku. Rozliším
elipsoidickou výšku GNSS od výšky nad mapovým výškovým systémem; pro
první analýzu sleduji zejména místní složku Up. Po přemístění základny
nebo novém survey-in nesmím bez dalšího překrýt stopy z různých dní.
Pro takové srovnání znovu ustavím společný referenční bod nebo doložím
transformaci mezi soustavami. Jde o můj návrh vyhodnocení, nikoli tvrzení
o již ověřené přesnosti.

\subsection{Proč simpleRTK2B místo hotového Holybro}
\label{sec:base-board}

Jako alternativu pro \emph{základnu}, nikoli pro rychlý rover, jsem
posuzoval Holybro H-RTK F9P Helical. Varianta UART nabízí USB-C a další
UART; provedení DroneCAN má rovněž vyvedený UART2. Výrobce uvádí
raw výstup \SI{20}{Hz} a RTK polohu \SI{8}{Hz}. Součástí sestavy je
helická anténa, kryt a kabeláž.\cite{holybro-f9p-uart,holybro-f9p-can}
Omezení rychlosti polohového výstupu samo o sobě tuto alternativu
nevylučuje jako statickou základnu: četnost korekcí není totožná
s požadovanou četností řešení pohybujícího se X5.

\begin{table}[htbp]
  \centering
  \small
  \caption{Úvaha o provedení základny; nejde o porovnání kvality GNSS čipu}
  \label{tab:base-board}
  \begin{tabularx}{\textwidth}{@{}L{0.25\textwidth}YY@{}}
    \toprule
    Varianta & Praktický přínos & Důsledek pro můj návrh \\
    \midrule
    Holybro F9P Helical, UART
      & hotové mechanické provedení a anténa
      & vhodná alternativa; menší volnost v uspořádání vlastní skříně \\
    Holybro F9P Helical, DroneCAN
      & integrace do existující sběrnice DroneCAN, také UART2
      & CAN nyní nevyužiji; není zdrojem vyšší přesnosti GNSS \\
    simpleRTK2B Budget
      & přístupné UART, USB, PPS a rozšiřující konektory
      & volím pro vlastní montáž a diagnostiku; anténa a kryt jsou zvlášť \\
    \bottomrule
  \end{tabularx}
\end{table}

Obě řešení využívají ZED-F9P, proto z názvu desky nevyvozuji rozdíl
v přesnosti přijímače. U simpleRTK2B oceňuji snadné připojení ESP32
a měřicích přístrojů, vlastní volbu antény a možnost objednat oba GNSS
přijímače společně.\cite{ardusimple-f9p-guide}
Ceny je nutné porovnávat včetně antény, kabeláže a krytu, nikoli jako
holou desku proti kompletnímu Holybro. CAN by dával smysl pro delší
robustní sběrnici nebo návaznost na autopilot; pro krátké propojení uvnitř
základny volím jednodušší UART.

\subsection{Samostatná základna s ESP32 a webovou správou}

Cílově chci provozovat základnu bez trvale připojeného notebooku.
ZED-F9P bude vytvářet RTCM, zatímco XIAO ESP32S3 obstará konfiguraci,
stavovou diagnostiku, webové rozhraní a obsluhu Wio-SX1262. GNSS připojím
obousměrně přes UART1. Podle revize desky ověřím nastavení IOREF na
\SI{3.3}{V}, společnou zem, napájení a nekolidující piny XIAO.
Sada XIAO/Wio není zásuvný XBee modul a do patice simpleRTK2B se přímo
nezasouvá.\cite{ardusimple-f9p-guide,seeed-xiao-wio}

ESP32 podporuje režim Wi-Fi přístupového bodu i klienta a může obsluhovat
webový server.\cite{esp32-wifi}
Z toho odvozuji možnost místní správy telefonem bez internetu. Vlastní
rozhraní má zobrazovat survey-in, použité souřadnice, stav záznamu,
RTCM tok a vytížení rádia; změnu konfigurace během jízdy omezím.
Vzdálená internetová správa je jiný režim: vyžaduje dostupnou síť,
například hotspot, a zabezpečený přístup. Zakoupený kit neobsahuje
mobilní modem ani hotovou aplikaci pro správu RTK základny.

Samostatnost nesmí být vykoupena ztrátou raw dat pro PPK. Do finální
základny proto musím doplnit a ověřit místní úložiště, například microSD
s vhodným rozhraním, a nezávislé napájení. Firmware oddělí ukládání
observací od odesílání RTCM; výpadek rádia nesmí zastavit zápis.
Rozhraní, kapacita a cena úložiště zatím nejsou uzavřené ani zahrnuté
v oceněném košíku. Samotná interní flash není předem ověřenou náhradou
dlouhého surového záznamu.

V první laboratorní etapě konfiguraci a raw log zajistí existující
týmový notebook přes USB. RTCM pak mohu porovnat přímou UART cestou
s hostitelským přenosem. RTKBase či ELT\_RTKBase zůstávají alternativou
pro základnu s linuxovým počítačem, ne firmwarem pro XIAO.
\cite{rtkbase,elt-rtkbase}
Notebook odpojím až po ověření místního loggeru a zotavení po výpadku
napájení. Druhé Raspberry Pi se kvůli správě základny zatím neobjednává.

\subsection{Volba čtyř sad XIAO ESP32S3 + Wio-SX1262}

Původně uvažované dva modemy Waveshare SX1262 LoRa HAT nahrazuji
\textbf{čtyřmi sadami Seeed Studio XIAO ESP32S3 + Wio-SX1262},
kód RPiShopu SEE319. Cena aktualizovaného seznamu je 339~Kč s DPH za sadu.
\cite{rpishop-xiao-wio}
Dvě sady tvoří aktivní spoj základna--rover. Pro zbývající dvě plánuji
laboratorní ověřování firmwaru a provozní rezervu. Neplánuji povinnou mesh
síť ani automatické přeposílání přes všechny čtyři jednotky.

Sada pro pásmo 862--930~MHz obsahuje XIAO, rádiovou desku a antény.
USB datový kabel není součástí. Předinstalovaný Meshtastic není v návrhu
zaměněn za transparentní RTCM modem: počítám s jeho nahrazením vlastním
firmwarem pro bod--bod přenos. Výrobce uvádí dodanou FPC anténu jako
testovací; pro spolehlivý provoz na trati ověřím vhodnější anténu a montáž.
\cite{seeed-xiao-wio}
\enquote{Nízkofrekvenční LoRa} zde znamená sub-GHz pásmo kolem 868~MHz,
nikoli fyzikální pásmo LF ani variantu 433~MHz.

Cílová datová cesta bude:
základna ZED-F9P, UART/XIAO vysílače, rádiový spoj,
XIAO přijímače, USB/Raspberry Pi a vstup RTCM mosaic-X5.
Přenosový rámec navrhnu s číslem sekvence, délkou, identifikací fragmentu
a kontrolou integrity. Neúplné nebo zastaralé zprávy zahodím a tuto událost
zaznamenám. Obnova spojení nesmí spustit přehrávání staré fronty korekcí.

\subsection{Propustnost, dosah a rádiová omezení}

Kilometrová trať je cílem ověřovacího testu, nikoli katalogovou garancí
zakoupené sady. Pro orientaci lze použít vztah pro volnoprostorový útlum
\begin{equation}
 L_\mathrm{FS}=32{,}44+
 20\log_{10}(f_\mathrm{MHz})+20\log_{10}(d_\mathrm{km}).
\end{equation}
Na \SI{1}{km} vychází přibližně \SI{91.2}{dB} pro 868~MHz a
\SI{100.0}{dB} pro 2,4~GHz. Rozdíl asi \SI{8.8}{dB} podporuje mou volbu
sub-GHz spoje, ale nezahrnuje antény, zastínění vozidlem, terén,
Fresnelovu zónu ani rušení.

RTCM základny nemusí běžet stejně rychle jako výstup roveru. Začnu
měřením skutečného datového toku při omezené sadě potřebných zpráv, například
MSM4 při \SI{1}{Hz}, a ověřím zachování FIX na X5. Nejde o předem potvrzenou
konečnou konfiguraci. Teprve z naměřených velikostí paketů zvolím spreading
factor, šířku pásma, kódování a případnou fragmentaci.

Pro každou zprávu spočtu skutečný čas ve vzduchu včetně preambule a režie.
Požadovanou dobu vysílání lze zapsat
\begin{equation}
 D_\mathrm{req}=\frac{\sum_i T_{\mathrm{air},i}}{T_\mathrm{pozorování}}.
\end{equation}
Ztráty a opakování tento podíl dále zvětšují. Nejvyšší přenosová rychlost
a nejlepší citlivost transceiveru nejsou dosažitelné současně při jednom
nastavení.\cite{sx1262}

Podle českého VO-R/10/05.2025-5 se limity liší mezi jednotlivými úseky
pásma. Například úsek 868,0--868,6~MHz pracuje s limitem 25~mW e.r.p.
a při režimu omezeném klíčovacím poměrem s 1\,\%; úsek 869,4--869,65~MHz
má odlišné podmínky včetně možnosti 10\,\%. Oprávnění stanovuje také
podmínky přístupu ke spektru a šířky pásma.\cite{ctu-vor10}
Tyto hodnoty nepředstavují hotovou konfiguraci modemu.

Pouhý název SX1262, detekce CAD či jednoduchý test obsazení kanálu
neprokazují vyhovující LBT režim. Při vlastním firmwaru musím doložit soulad
konkrétní konfigurace. Široký kanál nesmí přesahovat zvolený povolený
úsek; nelze například automaticky přenést rychlost pro 500kHz kanál do
úseku širokého pouze 250~kHz. Podmínky znovu ověřím před provozem a pro zemi
každého závodu.

Pokud RTCM nevychází do dovoleného vysílacího času a současně nesplní
požadované stáří korekcí, bude nutná úprava zpráv nebo přenosové cesty.
Dosah sám o sobě takový problém neřeší. Před jízdou je rozhodující
společný test propustnosti, rádiové shody a skutečné dostupnosti RTK FIX.

\subsection{NTRIP a chování při výpadku}

NTRIP umožňuje přenos RTCM přes IP síť a ponechávám jej jako kontrolní
a záložní cestu.\cite{bkg-ntrip}
RTK2Go a Centipede-RTK mohou sloužit při laboratorním srovnání podle
dostupnosti a podmínek služby; nejsou zaručenou součástí závodní
infrastruktury.\cite{rtk2go,centipede}

Ztráta rádiových korekcí nesmí zastavit logger. Systém uchová raw data
a označí změnu kvality GNSS. Současný raw log základny umožní pozdější
PPK i tehdy, když rádiový spoj během jízdy nefungoval. Pro offline PPK
tedy není nutné, aby korekce dorazily v reálném čase; pro online RTK je
jejich včasný příjem zásadní.

\section{Antény, mechanika a napájení}

\subsection{Od Budget Survey ke dvěma ANN-MB2-00}

Pro rover i základnu volím stejnou anténu \textbf{u-blox ANN-MB2-00},
celkem tedy dva kusy. Každá zahrnuje 5m kabel se SMA konektorem.
\cite{ardusimple-annmb2}
Původně jsem k jedné ANN-MB2 zvažoval Budget Survey Tripleband IP66 za
98~EUR bez DPH. Její praktickou výhodou je geodetický závit 5/8 palce
a odpojitelný kabel TNC--SMA o délce \SI{2.5}{m}.
\cite{ardusimple-budget-survey}
Po rozhodnutí vytvořit vlastní 3D tištěné kryty dávám přednost
sjednocení antén. Druhá ANN-MB2 stojí 61~EUR, takže výměna snižuje
oceněnou sestavu o 37~EUR, tedy 902,06~Kč při použitém kurzu.
Nejde o přidání třetí antény ani o dvouanténní heading.

ANN-MB2 má nízké provedení přibližně
$108{,}5\times92{,}9\times24{,}7$~mm a dokumentovaný referenční bod.
Její pevný 5m kabel však není ideální pro malou krabičku: při návrhu
musím vyřešit trasu, rezervu a odlehčení tahu, ne automaticky předpokládat
zkrácení kabelu.\cite{ublox-annmb2-ds}
Sjednocení mi usnadní montážní díly a případnou záměnu kusů. Není to
důkaz lepší RF kvality. Katalogové údaje o fázovém středu obou antén
nepovažuji bez shodné metodiky a montážních podmínek za přímý žebříček
přesnosti celé soustavy.

Třípásmová anténa navíc nezmění vybraný dvoupásmový ZED-F9P L1/L2 na
třípásmový přijímač. U základny využiji pouze společně podporované
signály; mosaic-X5 může využít širší příjem.\cite{ublox-f9p,septentrio-mosaic}
U obou desek před zapojením ověřím napájení aktivní antény a SMA připojení.

\subsection{Vlastní kryt a vodivá plocha pod anténou}

Anténu ponechám v jejím původním pouzdře a vlastní 3D tisk použiji jako
držák a kryt sestavy. Jako první variantu navrhuji nevodivé víko bez
uhlíkového či kovového plniva, anténu nad vodivou podložkou a elektroniku
na izolovaných distančních sloupcích pod ní. Je to konstrukční hypotéza
k ověření, nikoli záruka elektromagnetického stínění.

Výrobce uvádí charakterizaci fázového středu ANN-MB2 na kruhové vodivé
ploše o průměru \SI{120}{mm}; změna jejího tvaru a velikosti může změnit
výsledek.\cite{ublox-annmb2-ds}
Tento rozměr není hranicí, pod kterou anténa přestane fungovat, ani
zárukou přesnosti při jakékoli montáži. Pro prototyp mohu vyzkoušet
hliníkovou podložku podobné velikosti. Integrovaný magnet ji k hliníku
nepřidrží, proto navrhnu mechanické uchycení. Kryt, zemní plochu a
vzdálenost od LoRa, Wi-Fi a DC/DC ověřím srovnávacím měřením.

Sleduji C/N0, dostupnost FIX, opakovatelnost místní polohy a chování
při zapnuté elektronice; srovnám varianty s víkem a bez víka a různá
uspořádání podložky. Anténa musí mít volný výhled. Vlastní kryt nenahradí
kvalitní anténní konstrukci ani zkoušku vícecestného šíření.

\subsection{Umístění IMU a napájení}

IMU upevním na tuhý držák s definovanou orientací. Změřím třírozměrný
lever arm k fázovému středu GNSS antény a rotaci os vůči vozidlu
(boresight). Blízkost těžiště je vhodná pro inerciální měření, ale nové
využití magnetometru může vyžadovat kompromis kvůli magnetickému rušení.
Za riziko považuji také magnet v pouzdře ANN-MB2 a feromagnetické
montážní díly: IMU neumístím těsně pod anténní magnet bez ověření pole.
Případné pružné uložení připustím až po změření rezonancí a fázového posunu.

Napájecí návrh musí vycházet ze skutečné palubní sítě. Podle upřesnění
z 12.~října zajišťuje elektronické jištění palubní PMU. Návrh zahrne
ochranu proti přepólování a přepětí, filtraci a řízené ukončení záznamu při
poklesu napětí. DC/DC měnič a energetický doběh nejsou součástí původního
rozpočtu; napájení rozvíjí oddíl~\ref{sec:napajeni-vozidlo}. Ověřím také přístup ke GPIO pod M.2 HAT+, chlazení SSD
a Raspberry Pi a zajištění konektorů proti vibracím. Tyto části jsou
podmínkou funkční instalace, nikoli nepodstatným příslušenstvím.

\section{Výsledek výběru}

Pro první prototyp jsem zvolil mosaic-X5 na simpleRTK3B Pro, jednu přenosnou
základnu ZED-F9P, obě výše uvedené antény, Adafruit ICM-20948 a Raspberry Pi 5
se 2~GB RAM, Active Coolerem, M.2 HAT+ a vlastním SSD. Rádiový nákup obsahuje
čtyři XIAO/Wio sady. Elektronika je rozdělena pouze mezi RPiShop a ArduSimple.

Původní ISM330DHCX zůstává v rešerši jako vysvětlení postupu výběru a jako
technická alternativa. Novou IMU volím kvůli možnosti využít kompas a
praktické dostupnosti, současně přijímám potřebu ověřit její vyšší šum,
časování a magnetické rušení. Žádnou z těchto vlastností dosud nepovažuji
za experimentálně potvrzenou.

Hlavními podmínkami úspěchu zůstávají úplnost raw záznamů, prokazatelná
časová vazba a reprodukovatelná validace. Přibližně pěticentimetrový
rozestup inerciálních vzorků při \SI{100}{km\per\hour} sám neprokazuje
centimetrovou přesnost trajektorie.

Výběr podřizuji zejména relativní přesnosti vůči nehybné základně.
simpleRTK2B a společný typ antény mi dávají prostor pro vlastní mechaniku
a správu přes ESP32; nejde o tvrzení, že hotová sestava Holybro by byla
pro základnu nepoužitelná. Samostatný provoz bez notebooku je cílem
implementace, který vyžaduje místní raw logger, napájení a ověření.
